% Generated by support/export-artifacts.py; do not edit.
% Slug: ch10-loop
% Source: chapters/ch10.tex:133
% TikZ source: chapters/ch10.tex:134
% Original label: \label{fig:ch10-loop}
% Original caption (LaTeX source):
%   A counting loop whose head merges the first-entry value with values carried around the back edge, which forces the sign of \code{i} to bottom.
% Caption: A counting loop whose head merges the first-entry value with values carried around the back edge, which forces the sign of i to bottom.
% Setup: support/compiler-pseudocode.sty (TikZ libraries and shared styles).
\begin{tikzpicture}
\node[basic block] (b1) {B1: i = 0};
\node[basic block, below=1.1cm of b1] (b2) {B2: test i less 10};
\node[basic block, below left=1.1cm and 0.7cm of b2] (b3) {B3: i = i + 1};
\node[basic block, below right=1.1cm and 0.7cm of b2] (b4) {B4: use i};
\draw[cfg edge] (b1) -- (b2);
\draw[cfg edge] (b2) -- (b3);
\draw[cfg edge] (b2) -- (b4);
\draw[back edge] (b3) -- (b2);
\end{tikzpicture}

% Original figure LaTeX:
% \begin{figure}
% \begin{tikzpicture}
% \node[basic block] (b1) {B1: i = 0};
% \node[basic block, below=1.1cm of b1] (b2) {B2: test i less 10};
% \node[basic block, below left=1.1cm and 0.7cm of b2] (b3) {B3: i = i + 1};
% \node[basic block, below right=1.1cm and 0.7cm of b2] (b4) {B4: use i};
% \draw[cfg edge] (b1) -- (b2);
% \draw[cfg edge] (b2) -- (b3);
% \draw[cfg edge] (b2) -- (b4);
% \draw[back edge] (b3) -- (b2);
% \end{tikzpicture}
% \caption{A counting loop whose head merges the first-entry value with values carried around the back edge, which forces the sign of \code{i} to bottom.}
% \label{fig:ch10-loop}
% \end{figure}
